% =============================================================================
\section{Estruturas IIR e quantização de coeficientes}\label{sec:estr}
% =============================================================================

\subsection*{Teoria mínima}

\paragraph{Três estruturas equivalentes (mesma $H(z)$, comportamentos
diferentes em ponto fixo).}
\begin{itemize}
  \item \textbf{Forma Direta I:} implementa
        $H(z)=\frac{B(z)}{A(z)}$ com dois atrasos por polo; os coeficientes
        $a_k$ entram \emph{expandidos}
        $A(z)=1+a_1z^{-1}+\cdots+a_Nz^{-N}$.
  \item \textbf{Cascata (SOS):} $H(z)=\prod_i H_i(z)$ com cada $H_i$ de 2ª
        ordem (polos complexos conjugados juntos). Cada seção é quantizada
        \emph{independentemente}.
  \item \textbf{Paralela / grade (lattice):} soma de frações parciais /
        coeficientes de reflexão $k_m$ — cada $|k_m|<1$ garante estabilidade
        (usada no Desafio 07 e na codificação de voz).
\end{itemize}

\paragraph{Sensibilidade de raízes (por que a cascata salva).} Seja
$A(z)=\prod_i(z-p_i)$. Derivando $\sum_k a_kp_i^{-k}=0$ em relação a $a_k$:
\[
  \frac{\partial p_i}{\partial a_k}
  =-\frac{p_i^{\,N-k}}{\prod_{j\neq i}(p_i-p_j)} .
\]
O denominador contém \textbf{todos os pares de distâncias entre polos}. Logo:
(i) polos \emph{próximos entre si} (cascata de 6ª ordem em Forma Direta com
polos agrupados em 3 pares) têm sensibilidade enorme; (ii) a cascata separa o
produto em fatores de 2ª ordem, onde só \emph{um} par aparece no denominador —
sensibilidade de 2ª ordem, não de 6ª.

\paragraph{Quantização em ponto fixo.} Com $b$ bits na faixa
$[-L,+L]$, o passo é
\[
  q=\frac{2L}{2^{b}},\qquad
  \hat{a}_k=\left\lfloor\frac{a_k}{q}\right\rfloor q\ (\text{truncamento}),
\]
e a perturbação máxima por coeficiente é $q$. O erro relativo importa:
coeficientes pequenos (como $a_2^{2}$ de uma seção com $r=0{,}92$,
$0{,}846$) sofrem muito mais que coeficientes grandes.

\ideia{Em hardware de precisão fixa: \textbf{quantize sempre cada seção SOS
separadamente}. A diferença aparece de forma dramática quando $q$ cresce
(poucos bits) — veja a varredura do Exercício 14.2.}

\subsection*{Exercícios resolvidos}

% ---------------------------------------------------------------- 14.1
\exer{\textbf{(Desafio 06)} Um filtro de 6ª ordem tem polos
$r_ie^{\pm j\theta_i}$ com $(r,\theta)=(0{,}92,0{,}1\pi),\ (0{,}95,0{,}25\pi),\
(0{,}98,0{,}4\pi)$. Escreva o denominador na Forma Direta I e calcule
$a_1,\dots,a_6$.}

\textbf{Solução.} $A(z)=\prod_i(1-2r_i\cos\theta_i z^{-1}+r_i^2z^{-2})$.
As raízes complexas são
\[
  0{,}874972\pm j0{,}284296;\quad 0{,}671751\pm j0{,}671751;\quad
  0{,}302837\pm j0{,}932035 .
\]
Expandindo o produto (valores exatos obtidos numericamente):
\[
  \begin{array}{ll}
    a_1=-3{,}6991201832, & a_2=+6{,}9339730023,\\
    a_3=-8{,}1706448593, & a_4=+6{,}3467631101,\\
    a_5=-3{,}0715525712, & a_6=+0{,}7336265104 .
  \end{array}
\]
Note os coeficientes grandes e com sinais alternados: é exatamente esse
``espalhamento'' que amplifica o efeito da quantização.

\resp{$a_1=-3{,}69912$, $a_2=6{,}93397$, $a_3=-8{,}17064$, $a_4=6{,}34676$,
$a_5=-3{,}07155$, $a_6=0{,}73363$.}

% ---------------------------------------------------------------- 14.2
\exer{Quantize os coeficientes de 14.1 em ponto fixo com 12 bits na faixa
$[-8,+8]$ (truncamento), recalcule os polos e discuta a estabilidade. Repita
para 8 bits e explicite a diferença entre Forma Direta e Cascata.}

\textbf{Solução.} Passo de quantização:
\[
  q=\frac{2\cdot8}{2^{12}}=0{,}00390625 .
\]
Coeficientes quantizados:
\[
  \hat{a}_1..\hat{a}_6=\{-3{,}6953125;\ 6{,}93359375;\ -8{,}16796875;\
  6{,}34375;\ -3{,}0703125;\ 0{,}73046875\}.
\]
Polos resultantes:
\[
  |z|=\{0{,}97882626,\ 0{,}97882626,\ 0{,}96427740,\ 0{,}96427740,\
  0{,}90550994,\ 0{,}90550994\},
\]
todos $<1$: com \textbf{12 bits} a Forma Direta ainda é estável, mas o polo
projetado em $0{,}98$ caiu para $0{,}97883$ (variação de $1{,}2\times10^{-3}$).

A fragilidade aparece com menos bits. Varredura de $b$ bits (passo
$q=16/2^{b}$), mostrando $\max_i|p_i|$:
\begin{center}
\begin{tabular}{cccc}
\toprule
$b$ & $q$ & Forma Direta & Cascata (SOS)\\
\midrule
6  & 0{,}25000000 & $1{,}246715$ \textbf{instável} & $0{,}866025$ ok\\
7  & 0{,}12500000 & $1{,}107362$ \textbf{instável} & $0{,}935414$ ok\\
8  & 0{,}06250000 & $1{,}002093$ \textbf{instável} & $0{,}968246$ ok\\
9  & 0{,}03125000 & $0{,}997348$ ok & $0{,}968246$ ok\\
10 & 0{,}01562500 & $0{,}988678$ ok & $0{,}976281$ ok\\
12 & 0{,}00390625 & $0{,}978826$ ok & $0{,}978280$ ok\\
\bottomrule
\end{tabular}
\end{center}
Resumo quantitativo: a Forma Direta é estável para todo $b\ge9$ bits (em
$8$ bits já estoura, $\max|p|=1{,}002093$); a Cascata é estável para todo
$b\ge6$ bits.

\resp{Com 12 bits: Forma Direta estável ($\max|p|=0{,}978826$) e Cascata
estável ($0{,}978280$) — ambas servem. Com 8 bits: Forma Direta
\textbf{instável} ($\max|p|=1{,}0021$), Cascata ainda estável ($0{,}9682$).
Ou seja, o ganho da estrutura em cascata é de \textbf{2--3 bits de precisão}.}

% ---------------------------------------------------------------- 14.3
\exer{Quantize as três seções SOS com os mesmos 12 bits e verifique os polos
de cada seção.}

\textbf{Solução.} Cada seção tem
$H_i(z)=\dfrac{1}{1-2r_i\cos\theta_iz^{-1}+r_i^2z^{-2}}$:
\begin{center}
\begin{tabular}{cccccc}
\toprule
$r$ & $\theta$ & $a_1$ & $\hat{a}_1$ & $a_2$ & $\hat{a}_2$\\
\midrule
$0{,}92$ & $0{,}10\pi$ & $-1{,}7499439900$ & $-1{,}7460937500$ & $0{,}8464000000$ & $0{,}8437500000$\\
$0{,}95$ & $0{,}25\pi$ & $-1{,}3435028843$ & $-1{,}3398437500$ & $0{,}9025000000$ & $0{,}9023437500$\\
$0{,}98$ & $0{,}40\pi$ & $-0{,}6056733090$ & $-0{,}6054687500$ & $0{,}9604000000$ & $0{,}9570312500$\\
\bottomrule
\end{tabular}
\end{center}
Polos de cada seção quantizada:
\[
  |z|=\{\,0{,}91855865,\ 0{,}94991776,\ 0{,}97827974\,\},
\]
todos dentro do círculo unitário. A seção mais crítica (projeto $r=0{,}98$)
desloca o polo para $0{,}978280$ — mesmo assim estável.

\resp{$\max|p_i|=0{,}97827974<1$: a cascata com 12 bits mantém a
estabilidade e o erro de posição dos polos é da ordem de $q=0{,}0039$ em cada
seção, \emph{não} amplificado pela ordem do filtro.}
